% ============================================================
% Lembar Persetujuan Pengajuan Magang Semester 7
% Standalone (article, Times) — tidak pakai shared/preamble.tex
% Build: make pdf  → output/main.pdf
% ============================================================
\documentclass[12pt]{article}

\usepackage[a4paper, margin=1in]{geometry}
\usepackage{graphicx}
\usepackage{tabularx}
\usepackage{array}
\usepackage{setspace}
\usepackage[T1]{fontenc}
\usepackage{mathptmx} % Times New Roman, sama dengan projek lain di repo ini

\pagestyle{empty}

% Kolom seragam untuk blok identitas: label lebar tetap (titik dua sejajar),
% isi rata kiri yang boleh turun baris, NIP/NRP kolom natural (tidak pernah pecah).
\newcolumntype{Y}{>{\raggedright\arraybackslash}X}
% Lebar kolom NIP/NRP dikunci = lebar string NIP terpanjang, supaya NIP dan NRP
% mulai di titik yang sama dan tidak pernah pecah baris.
\newlength{\KolomNip}


% ================= DATA =================
\newcommand{\NamaPembimbingSatu}{Prof. M. Udin Harun Al Rasyid, S.Kom, Ph.D}
\newcommand{\NipPembimbingSatu}{198108082005011001}
\newcommand{\NamaPembimbingDua}{Grezio Arifiyan Primajaya, S.Kom, M.Kom}
\newcommand{\NipPembimbingDua}{199208262022031007}

\newcommand{\NamaMahasiswa}{Muhammad Rafi Rizaldi}
\newcommand{\NrpMahasiswa}{3123600001}
\newcommand{\Jurusan}{Teknik Informatika}
\newcommand{\JudulProyekAkhir}{Implementasi Sistem Monitoring Detak Jantung Kontinu Berbasis IoT Edge Computing untuk Deteksi Dini Aritmia}

\newcommand{\NamaPerusahaan}{PT Aigra Eon Indonesia}
\newcommand{\TanggalMulai}{3 Agustus 2026}
\newcommand{\TanggalSelesai}{3 Desember 2026}
\newcommand{\TanggalSurat}{30 Juli 2026}

\newcommand{\LogoPens}{../proposal-pa/gambar/logo-pens.png}

% ================= DOKUMEN =================
\begin{document}

% diukur di dalam document supaya font (Times 12pt) sudah aktif
\settowidth{\KolomNip}{NIP. 199208262022031007}

% ---------- Kop / header berbingkai ----------
\noindent
% tabularx: kolom teks (X) memakan sisa lebar, jadi kop mentok margin kiri-kanan
\begin{tabularx}{\textwidth}{|c|>{\centering\arraybackslash}X|}
\hline
\raisebox{-0.5\height}{%
  \IfFileExists{\LogoPens}%
    {\includegraphics[width=2.2cm]{\LogoPens}}%
    {\parbox[c][2.4cm][c]{2.2cm}{\centering LOGO\\PENS}}%
} &
\textbf{Politeknik Elektronika Negeri Surabaya}\par
\vspace{2pt}\hrule\vspace{6pt}
\textbf{LEMBAR PERSETUJUAN}\par
\textbf{PENGAJUAN MAGANG SEMESTER 7}
\\ \hline
\end{tabularx}

\vspace{1.5cm}

% ---------- Judul ----------
\begin{center}
  {\large\bfseries LEMBAR PERSETUJUAN\\
  PENGAJUAN MAGANG SEMESTER 7}
\end{center}

\vspace{0.8cm}

\noindent Yang bertanda tangan di bawah ini :

\vspace{0.4cm}

\noindent
\begin{tabularx}{\textwidth}{@{}p{4.2cm}@{ : }Y@{\hspace{10pt}}p{\KolomNip}@{}}
Nama Pembimbing 1 & \NamaPembimbingSatu & NIP. \NipPembimbingSatu \\[4pt]
Nama Pembimbing 2 & \NamaPembimbingDua  & NIP. \NipPembimbingDua  \\
\end{tabularx}

\vspace{0.3cm}

\noindent Sebagai Dosen Pembimbing dari mahasiswa :

\vspace{0.3cm}

\noindent
\begin{tabularx}{\textwidth}{@{}p{4.2cm}@{ : }Y@{\hspace{10pt}}p{\KolomNip}@{}}
Nama    & \NamaMahasiswa & NRP. \NrpMahasiswa \\[4pt]
Jurusan & \Jurusan       & \\
\end{tabularx}

\vspace{4pt}

\noindent
\begin{tabularx}{\textwidth}{@{}p{4.2cm}@{ : }Y@{}}
Judul Proyek Akhir & \JudulProyekAkhir \\
\end{tabularx}

\vspace{0.5cm}

\noindent Dengan ini menyatakan :

\vspace{0.2cm}

\noindent
Menyetujui mahasiswa tersebut di atas untuk mengajukan magang di perusahaan
\textbf{\NamaPerusahaan} untuk periode magang sejak tanggal \TanggalMulai{}
sampai tanggal \TanggalSelesai.

\vspace{0.3cm}

\noindent Demikian untuk diketahui.

\vspace{1cm}

% ---------- Tanggal & tanda tangan ----------
\begin{center}
Surabaya, \TanggalSurat\\
Mengetahui / Menyetujui,
\end{center}

\vspace{0.5cm}

\noindent
\begin{tabularx}{\textwidth}{@{}>{\centering\arraybackslash}X >{\centering\arraybackslash}X@{}}
Dosen Pembimbing 1 & Dosen Pembimbing 2 \\[2.2cm]
(\NamaPembimbingSatu) & (\NamaPembimbingDua) \\
{\small NIP. \NipPembimbingSatu} & {\small NIP. \NipPembimbingDua} \\
\end{tabularx}

\end{document}
